% =============================================================================
% Chapter 6: Conclusions
% =============================================================================

\chapter{Conclusions}
\label{ch:conclusions}

This chapter summarizes the work presented in this thesis, restates the contributions, discusses the limitations of the current implementation, and outlines directions for future research and development.


% =============================================================================
\section{Summary}
\label{sec:summary}

Multi-cluster Kubernetes deployments are increasingly common, yet they suffer from resource fragmentation: some clusters have idle capacity while others face resource pressure. This thesis presented a system that addresses this imbalance through centralized resource brokerage with automatic cluster peering.

The system consists of two components: a \textbf{broker} that collects resource advertisements from all participating clusters, processes reservation requests through a scoring-based decision engine, and manages the reservation lifecycle; and an \textbf{agent} deployed in each cluster that monitors local resources, publishes advertisements, receives instructions, and automatically establishes Liqo peering with the assigned provider.

The broker uses a \texttt{Reserved} field mechanism to prevent race conditions under concurrent reservations: when a provider is selected, its available capacity is atomically decremented before the reservation is confirmed. This ensures that concurrent requests see accurate availability and cannot double-book the same resources.

All communication is secured with mutual TLS (mTLS), where cluster identity is derived from the client certificate's Common Name. This model eliminates the need for separate authentication tokens or shared identity providers, making the system suitable for deployments where clusters belong to different teams or organizations.

The experimental evaluation, conducted on a 96-core server with Kind clusters, demonstrated:

\begin{itemize}
    \item Broker CPU scales linearly at 0.23\% per agent, with constant $\sim$40~MB memory.
    \item Reservation decisions complete in under 200~ms.
    \item Up to 10 concurrent reservations are handled with zero timeouts and zero double-bookings.
    \item The decision engine achieves 83--100\% accuracy in selecting the optimal provider.
    \item Advertisement propagation delay stabilizes at 1.4 seconds.
    \item A single agent consumes only 41.25~MB of memory and negligible CPU.
\end{itemize}


% =============================================================================
\section{Contributions}
\label{sec:final-contributions}

The specific contributions of this thesis are:

\begin{enumerate}
    \item \textbf{Centralized broker architecture} for multi-cluster Kubernetes resource sharing, implemented as a Kubernetes operator with a REST API. The broker collects real-time resource metrics, manages the reservation lifecycle, and serves as the single source of truth for cluster availability.

    \item \textbf{Reserved field mechanism} for race-condition prevention. The broker atomically updates a \texttt{Reserved} field in the provider's \texttt{ClusterAdvertisement} upon each reservation, and the agent preserves this field during advertisement updates through a two-step read-then-write protocol.

    \item \textbf{Scoring-based decision engine} that selects the optimal provider by evaluating projected post-reservation utilization. The algorithm filters candidates by resource sufficiency and active status, then ranks them by remaining headroom, distributing load across the federation.

    \item \textbf{Certificate-based identity model} using mTLS with cert-manager. Cluster identity equals the certificate Common Name, providing mutual authentication without additional infrastructure.

    \item \textbf{Automatic Liqo peering integration.} The agent's \texttt{ReservationInstructionReconciler} automatically triggers \texttt{liqoctl peer} when an instruction is received, creating virtual nodes and encrypted WireGuard tunnels without manual intervention.

    \item \textbf{Comprehensive evaluation suite} of eight automated tests covering scalability, latency, concurrency, accuracy, resource tracking, and freshness, all with reproducible scripts and published results.

    \item \textbf{Open-source implementation} in Go, using standard Kubernetes tooling (controller-runtime, Kubebuilder, CRDs), available for both academic research and practical deployment.
\end{enumerate}


% =============================================================================
\section{Limitations}
\label{sec:limitations-final}

The current implementation has the following limitations:

\begin{enumerate}
    \item \textbf{Single broker (no high availability).} The broker is a single point of failure. While controller-runtime supports leader election for multi-replica deployments, the HTTP API server and decision engine have not been tested in an active-passive configuration. A broker failure would prevent new reservations until the process restarts, though existing peering relationships would be unaffected.

    \item \textbf{Polling-based instruction delivery.} Agents discover new reservations by polling the broker at a configurable interval (default: 30 seconds). This introduces latency between the broker's decision and the agent's action. A push-based mechanism (WebSocket, gRPC streaming, or MQTT) would reduce this to near-real-time.

    \item \textbf{Single-resource scoring.} The decision algorithm considers only CPU and memory availability with equal weight. It does not account for network latency between clusters, energy cost, geographic proximity, or specialized hardware (GPU models, FPGA). The scoring formula uses fixed weights rather than configurable policies.

    \item \textbf{No reservation cancellation or modification.} Once a reservation is created, it can only expire or be deleted. There is no mechanism to resize a reservation (e.g., increase CPU from 500m to 1 core) or transfer it to a different provider without creating a new reservation.

    \item \textbf{Single-machine evaluation.} All experiments were conducted with Kind clusters on a single server. This environment does not reflect the network latency, bandwidth constraints, and heterogeneous hardware of a real multi-site deployment.

    \item \textbf{Static kubeconfig-based peering.} The Liqo integration requires pre-distributed kubeconfig files for all clusters. In a dynamic federation where clusters join and leave, a more flexible discovery and authentication mechanism would be needed.
\end{enumerate}


% =============================================================================
\section{Future Work}
\label{sec:future-work}

Several directions for extending the system are identified:

\begin{enumerate}
    \item \textbf{High availability and broker federation.}
    Deploy the broker in a highly available configuration with leader election and state replication. For large-scale federations, multiple brokers could be organized hierarchically (e.g., regional brokers coordinating with a global broker), or in a peer-to-peer topology where brokers exchange advertisements.

    \item \textbf{Push-based notification.}
    Replace the polling-based instruction delivery with WebSocket or gRPC streaming connections. This would reduce the end-to-end latency from the current $\sim$53 seconds (dominated by polling interval) to under 1 second, making the system suitable for latency-sensitive use cases.

    \item \textbf{Advanced scheduling policies.}
    Extend the scoring algorithm to incorporate multiple objectives: network latency between requester and provider, energy efficiency (prefer clusters powered by renewable energy), cost optimization (prefer cheaper providers), data locality (prefer clusters in the same region as the data), and hardware affinity (match GPU model requirements).

    \item \textbf{Multi-resource support.}
    Extend the data model and decision engine to handle GPU allocation (with model-specific matching), persistent storage (with capacity and IOPS requirements), and network bandwidth reservations.

    \item \textbf{Dynamic Liqo peering.}
    Replace the static kubeconfig-based peering with a dynamic mechanism where the broker facilitates mutual authentication between clusters at peering time, potentially using short-lived tokens or delegated certificates.

    \item \textbf{Reservation lifecycle enhancements.}
    Add support for reservation modification (scaling up or down), migration (moving a reservation to a different provider), and preemption (higher-priority reservations displacing lower-priority ones).

    \item \textbf{Observability and monitoring.}
    Integrate Prometheus metrics and Grafana dashboards for real-time visibility into broker load, reservation throughput, cluster utilization, and peering status across the federation.

    \item \textbf{Production deployment.}
    Package the broker and agent as Helm charts for standard Kubernetes deployment, with support for Kubernetes-native scaling, rolling updates, and integration with existing cert-manager installations.
\end{enumerate}

\bigskip

\noindent In conclusion, this thesis demonstrates that centralized resource brokerage combined with automatic cluster peering is a viable and effective approach to multi-cluster Kubernetes resource sharing. The system achieves sub-200~ms decision latency, handles concurrent requests without double-booking, and bridges the gap between resource allocation and actual cluster interconnection through Liqo integration. The open-source implementation provides a foundation for both further research and practical deployment in multi-cluster environments.
